\chapter{Conclusion and Future Work}
\label{ConclusionAndFutureWork}

\section{Conclusion}

This thesis presented WATCHOUT ISEQL, a forensic multimodal surveillance framework that combines visual perception from a VLM with object re-identification and environmental sound analysis from a LALM through a shared interval-based temporal representation. The study was organised as a three-condition ablation: visual-only (A), audio-only (B), and multimodal union (C = A $\cup$ B).

The evaluation on a 30-scene dataset across four VLM models and two audio providers yields three principal findings:

\begin{enumerate}
  \item \textbf{Object re-identification is essential.} Without it, all four VLM models produce zero true positives. The two-phase tracking prompt with block-overlap matching fixes the zero-event bug identified in prior work.
  \item \textbf{The LALM achieves perfect precision.} Qwen2-Audio-7B-Instruct attains precision of 1.000 at the optimal 2.5 s / 1.25 s configuration, with an F1 of 0.750 ,  substantially outperforming PANNs CNN14 (F1 = 0.429).
  \item \textbf{Multimodal union outperforms both unimodal baselines.} The best combination (Gemini 3.6 Flash + Qwen2-Audio) achieves 29/30 TP (F1 = 0.892, 6 FP, 1 FN), and the claim $C \ge \max(A, B)$ holds on all 30 scenes.
\end{enumerate}

The union approach, concatenating independent modality detections without temporal cross-modal joins ,  is shown to be more effective than fusion approaches that rely on temporal alignment, which reduce recall. Every detection is the output of a SQL query, providing a fully auditable pipeline suitable for forensic applications.

\section{Future Work}

\textbf{Real CCTV evaluation.} The present dataset is synthetic. Applying the same three-condition methodology to operational CCTV footage would test the generalisability of the findings. The evaluation harness supports adding new ground truth files without code changes.

\textbf{Stronger audio models.} Qwen2-Audio-7B is evaluated in this thesis. Stronger models may improve recall on the events where the current model fails (fight detection, 0/5). A head-to-head comparison of multiple LALMs on this surveillance benchmark would be valuable.

\textbf{Denser visual sampling.} The current 1.4 Hz frame rate is too sparse for fast events. A YOLO-based pre-filter \citep{redmon2016yolo} could identify actor bounding boxes and trigger denser VLM sampling on those regions, improving temporal resolution without increasing the total number of VLM calls.

\textbf{Online parameter tuning.} The per-query delta parameters and confidence thresholds are static. An online loop that adjusts these based on operator feedback would improve adaptability to different environments.

\textbf{Temporal IoU metrics.} The current evaluation uses binary verdicts. Incorporating temporal IoU would measure how precisely the detected interval aligns with the ground truth, providing a finer-grained view of system performance.

\textbf{Multi-camera support.} The interval representation extends naturally to multiple cameras and PTZ configurations, but the VLM prompt must handle changing viewpoints and the temporal join logic must account for inter-camera alignment.
